\section{Interpolation at logarithmic centers}\label{sec:interpolation}

Let \(m,K\ge1\), and fix positive rational \(w_0,v_0,\theta\)
with \(0<\theta<1\). For further positive rational weights \(w_i\),
set
\[
 W=(w_0,w_1,\ldots,w_m),\qquad
 V=(v_0,v_1,\ldots,v_m),\qquad v_i=w_i/\theta\quad(i>0).
\]
Let \(\PW(H)\) be the space spanned by \(Y^hX^\alpha\) with
\(w_0h+w\alpha\le H\). Let \(\JV(H)\) be the quotient of
\(\C[[t,u_1,\ldots,u_m]]\) retaining precisely the coefficients
\[
 v_0\ell+\sum_i v_i\beta_i<H.
\]
The inequality for rows is strict, whereas the inequality for columns
is weak.

\begin{theorem}[Separated logarithmic interpolation]\label{thm:interpolation}
Assume
\begin{equation}\label{eq:volume}
 K(w_0/v_0)\theta^m<1.
\end{equation}
There are successive lower thresholds for \(w_1,\ldots,w_m\),
independent of the center coordinates, with this property.
For nonzero \(y_j\in\C\), \(0\le j<K\), and \(c_{ji}\in\C\),
the map
\begin{equation}\label{eq:jetmap}
 \begin{split}
 \PW(H)&\longrightarrow\bigoplus_{j<K}\JV(H),\\
 P&\longmapsto
 \left(P\bigl(y_j(1+t),
       (c_{ji}+u_i+\log(1+t))_{i=1}^m\bigr)\right)_{j<K}
 \end{split}
\end{equation}
is surjective for all sufficiently large \(H\) divisible by a fixed
positive integer, in either of the following cases:
\begin{enumerate}[label=(\roman*)]
\item the \(y_j\) are pairwise distinct;
\item all \(y_j=1\), each family \(c_{0i},\ldots,c_{K-1,i}\)
is pairwise distinct, and \(K\theta^m<1\).
\end{enumerate}
The divisibility requirement depends only on the weights. The lower
threshold for \(H\) may depend on every fixed datum, including the
centers.
\end{theorem}

Case (ii) is the common-fiber interpolation of \cite[Section 2]{OAIpi}.
We give the argument with both cases present to exhibit the exact
change required for case (i). In particular, a count of coefficients
alone is not used to assert surjectivity.

\subsection{The uniform local input and separated weights}

We use the weighted transverse multiplicity comparison of
\cite[Lemma 2.2]{OAIpi}. In the notation needed here it says the following.

\begin{lemma}[Weighted multiplicity comparison]\label{lem:multiplicity}
Give \(D\) affine coordinates positive degree weights \(\rho_a\).
Let commuting analytic vector fields form a frame on an open set,
with positive derivative costs \(\kappa_b\).
Suppose polynomials \(f_\ell\) have weighted degree at most \(N\),
and \(Z\) is a codimension-\(k\) component of their common zero set.
If every further frame derivative of cost less than
\(\epsilon N\) of every \(f_\ell\) vanishes on \(Z\), then, for
any coordinate normal basis \(A\) and frame normal basis \(B\),
\begin{equation}\label{eq:local-comparison}
 \prod_{a\in A}\rho_a
 \le 2k!\epsilon^{-k}\prod_{b\in B}\kappa_b.
\end{equation}
The constant is independent of the equations and of the component.
\end{lemma}
\begin{proof}
Work at a general smooth point of \(Z\) where both normal-basis
conditions hold and no other component passes through the point.
The commuting flows in the \(B\)-directions give transverse analytic
coordinates. Vanishing of the derivatives says that the equations
have no transverse terms of weight below \(\epsilon N\).
Restrict to the coordinate slice in the \(A\)-directions.
That slice is transverse, its common zero is isolated, and substitution
of the remaining analytic coordinates cannot introduce a transverse
term of smaller weight.
Thus the local quotient length is at least
\[
 \#\{\alpha\in\N^k:\textstyle\sum_{b\in B}\kappa_b\alpha_b
                  <\epsilon N\}
 \ge \frac{(\epsilon N)^k}{2k!\prod_{b\in B}\kappa_b}.
\]
For the last inequality, unit cubes indexed by the integer parts of
points of the open weighted simplex cover that simplex; even the
factor \(2\) is unnecessary for this count.

Generic \(k\) constant linear combinations of the restricted equations
still cut out an isolated zero. Their quotient length bounds the
preceding length from above, and weighted B\'ezout bounds it by
\(N^k/\prod_{a\in A}\rho_a\).
For clarity, this weighted version follows by choosing an integer \(M_0\)
with \(M_0\rho_a\in\N_{>0}\) and substituting
\(x_a=\lambda_a+z_a^{M_0\rho_a}\). Choose \(\lambda_a\) so that
the preimages of the point are unramified. There are
\(\prod_{a\in A}M_0\rho_a\) such preimages, each preserving the
local length, while the total degrees are at most \(M_0N\).
Ordinary B\'ezout for isolated zeros gives the asserted upper bound;
see \cite[Corollary VIII.3]{Mondal2021}.
Comparison of the two lengths proves \eqref{eq:local-comparison}.
\end{proof}

Choose \(\sigma_0\in\Q_{>0}\) small enough that
\begin{equation}\label{eq:sigma}
 \begin{split}
 (1+3\sigma_0)^{m+1}K(w_0/v_0)\theta^m&<1,\\
 (1+\sigma_0)\theta&<1,
 \end{split}
\end{equation}
and, in case (ii), also
\((1+3\sigma_0)^mK\theta^m<1\).
Put \(L_0=m+2\) and take the deliberately enlarged constant
\[
 C_*=
 2(m+2)^{m+2}\left(1+\frac{m+2}{\sigma_0}\right)^{m+2}.
\]
It dominates \(2k!(L_0/\sigma_0)^k\) for \(1\le k\le m\).
Choose the positive weights successively so that
\begin{equation}\label{eq:weight-separation}
 \prod_{a\in A}w_a>C_*\prod_{b\in B}v_b
\end{equation}
whenever \(A,B\subseteq\{0,\ldots,m\}\) have equal size between
\(1\) and \(m\), and the largest positive index in their symmetric
difference belongs to \(A\).
At that index the ratio in \eqref{eq:weight-separation} contains the
new weight linearly. Larger common indices contribute only
\(w_i/v_i=\theta\); all remaining factors were already fixed.
Finitely many inequalities therefore impose a finite lower threshold.
No point coordinate enters these thresholds.

The formal implementation can use a rectangular transverse-length
bound, with constant \(k^k(L_0/\sigma_0)^k\), at a sufficiently large
auxiliary degree. The displayed \(C_*\) also dominates that constant.
The auxiliary-degree threshold is independent of the varying component,
so either local formulation gives the same separated-weight argument.

\subsection{The curve inequality}

On the smooth complete model of an irreducible curve \(C\) meeting
the affine chart, define
\begin{equation}\label{eq:weighted-degree}
 \deg_W C=\sum_P
  \max\left\{0,-\frac{\ord_PY}{w_0},
          -\frac{\ord_PX_1}{w_1},\ldots,
          -\frac{\ord_PX_m}{w_m}\right\}.
\end{equation}
An identically zero coordinate contributes no pole.
At a branch through \(a_j=(y_j,c_{j1},\ldots,c_{jm})\), let
\begin{equation}\label{eq:contact}
 h_P=\min\left\{
  \frac{\ord_P(Y/y_j-1)}{v_0},
  \frac{\ord_P(X_i-c_{ji}-\log(Y/y_j))}{v_i}:1\le i\le m
 \right\}.
\end{equation}
The logarithm is a formal series at the branch, not a global choice
of a complex logarithm. These contacts are positive and finite on
a nonconstant branch.

\begin{proposition}\label{prop:curve}
For the separated weights above, in either case of
Theorem~\ref{thm:interpolation},
\begin{equation}\label{eq:curve-bound}
 \deg_W C\ge(1+\sigma_0)\sum_{P\mapsto a_j}h_P.
\end{equation}
\end{proposition}
\begin{proof}
Suppose the inequality fails. A count of coefficients and conditions
gives, for large auxiliary integers \(N\), a nonzero polynomial \(F_N\)
of \(W\)-degree at most \(N\), with local \(V\)-order at least
\((1+3\sigma_0)N\) at every center. The leading terms in the two counts
are
\[
 \frac{N^{m+1}}{(m+1)!\prod_aw_a},
 \qquad
 \frac{K(1+3\sigma_0)^{m+1}N^{m+1}}
      {(m+1)!\prod_av_a},
\]
and \eqref{eq:sigma} makes the first larger. Independence of the
conditions is not needed to produce \(F_N\).

On \(Y\ne0\), use the commuting frame
\[
 D_0=Y\partial_Y+\sum_{i=1}^m\partial_{X_i},
 \qquad D_i=\partial_{X_i}\quad(i>0).
\]
In the coordinates of \eqref{eq:jetmap}, these are
\((1+t)\partial_t,\partial_{u_i}\).
Each preserves the bound on polynomial degree; differentiating with
cost \(v_a\) lowers local weighted order by at most \(v_a\).
Thus every derivative of \(F_N\) of total cost at most
\(\sigma_0N\) has branch order at least
\((1+2\sigma_0)Nh_P\).
Its total pole order on the complete curve is at most \(N\deg_W C\).
The assumed failure of \eqref{eq:curve-bound} makes the former
total strictly larger. The zero--pole equality forces all these
derivatives to vanish identically on \(C\).

Set \(\delta_N=\sigma_0N/L_0\). Consider successively the common
zero sets of derivatives of cost at most \(r\delta_N\),
\(0\le r\le L_0\), on \(Y\ne0\).
Choose nested irreducible components containing the affine part of
\(C\), working backwards from the last level.
Their dimensions lie between \(1\) and \(m\); hence two adjacent
components coincide. Call this persistent component \(Z\), of
codimension \(k\), \(1\le k\le m\).
Further derivatives of cost less than \(\delta_N\) of its defining
equations from the earlier level vanish on \(Z\).
Lemma~\ref{lem:multiplicity} gives
\[
 \prod_{a\in A}w_a\le C_*\prod_{b\in B}v_b
\]
for every coordinate normal basis \(A\) and frame normal basis \(B\).
This argument is an instance of the persistent derivative-ideal method;
see also \cite[Section 5]{Philippon1986}.

If some positive coordinate direction is nontangent to \(Z\),
choose its largest index \(i\). Some coordinate normal basis contains
\(i\). Separation forces every frame normal basis to contain \(D_i\).
The span of the remaining frame vectors is
\(\ker(dX_i-dY/Y)\). Its failure to span the normal quotient implies
\[
 dX_i=dY/Y\quad\text{on }Z.
\]
This forces \(Y\) to be constant. Otherwise restrict to an algebraic
curve on which \(Y\) is nonconstant and pass to its complete model.
At a zero or pole of \(Y\), the residue of \(dY/Y\) is the
nonzero integer \(\ord_PY\), whereas the exact differential \(dX_i\)
has residue zero. This is impossible.
If every positive coordinate direction is tangent, properness of \(Z\)
forces its tangent space to be their span, so \(dY=0\) and the same
conclusion holds. In particular \(Y\) is constant on \(C\).

In case (i), \(C\) meets at most one prescribed center \(a_j\).
Choose \(X_\ell\) nonconstant on \(C\). The function
\(X_\ell-c_{j\ell}\) has total pole degree at most
\(w_\ell\deg_W C\), and has zero order at least \(v_\ell h_P\)
at every relevant branch. Consequently
\[
 \deg_W C\ge\theta^{-1}\sum_{P\mapsto a_j}h_P
          >(1+\sigma_0)\sum_{P\mapsto a_j}h_P,
\]
contradicting the supposition.

In case (ii), the constant is \(Y=1\).
In that fiber a second coefficient count is available because
\(K(1+3\sigma_0)^m\theta^m<1\).
It gives a nonzero polynomial in \(X_1,\ldots,X_m\) of degree
at most \(N\) and order at least \((1+3\sigma_0)N\) at all centers.
The same zero--pole argument makes its low-cost ordinary partial
derivatives vanish on \(C\).
For \(m=1\), a nonzero polynomial cannot vanish on this curve.
For \(m\ge2\), repeat the persistent-component argument in the fiber.
The coordinate and derivative frames now agree.
Two different normal-basis index sets would violate
\eqref{eq:weight-separation}, so there is just one.
Choose an index \(i\) in it. The other coordinate directions fail
to span the normal quotient, so \(dX_i=0\) on the component.
That coordinate is constant in characteristic zero.
Coordinatewise distinctness implies that the component, and hence
\(C\), meets at most one center. The preceding single-center
zero--pole argument again contradicts the supposed excess.
\end{proof}

\subsection{From curve degrees to the actual coefficient map}

We finish the proof of Theorem~\ref{thm:interpolation}; this is the
curve-to-jet passage of \cite[Section 2.3]{OAIpi}, with normalized
coordinates at arbitrary \(y_j\ne0\).
Choose an integer \(R\) for which every \(R/w_a\) and \(R/v_a\)
is a positive integer. Embed the affine space by all monomials of
\(W\)-degree at most \(R\), including \(1\), and take its projective
closure \(X\), with \(L=\OO_X(1)\).
Pure coordinate powers attain the maximum pole at each point of a
complete curve, so
\begin{equation}\label{eq:projective-degree}
 \deg(L|_C)=R\deg_W C.
\end{equation}

Replace each logarithm by a Taylor polynomial \(G_i(t)\) whose first
omitted term has \(t\)-weight greater than \(v_i\). The coordinates
\[
 t=Y/y_j-1,\qquad u_i'=X_i-c_{ji}-G_i(t)
\]
are algebraic regular coordinates at \(a_j\).
The replacement preserves \eqref{eq:contact} and induces inverse
filtration-preserving triangular changes on every finite coefficient
quotient. Define the finite-colength ideal
\[
 I_j=(t^{R/v_0},(u_1')^{R/v_1},\ldots,(u_m')^{R/v_m}).
\]
Patch these ideals with the unit ideal away from the centers.
Its branch order is \(Rh_P\).

Let \(p:X'\to X\) be the ordinary blow-up of this ideal \(I\), and
write \(I\OO_{X'}=\OO_{X'}(-E)\), \(A=p^*L\).
The curve inequality and \eqref{eq:projective-degree} show that
\(A-(1+\sigma_0)E\) is nef.
For a curve contracted by \(p\), this follows instead from relative
ampleness of \(-E\); curves avoiding the affine chart do not meet \(E\).
For large \(a>1\), \(aA-E\) is ample.
The identity
\[
 A-E=
 \frac{a-1}{a(1+\sigma_0)-1}\bigl(A-(1+\sigma_0)E\bigr)
 +\frac{\sigma_0}{a(1+\sigma_0)-1}(aA-E)
\]
therefore makes \(A-E\) ample; see
\cite[Corollary 1.4.10]{Lazarsfeld2004}.
These statements do not require a smooth or normal blow-up.

The large-degree graded-piece and vanishing theorems for the Rees
algebra give
\[
 p_*\OO_{X'}(-nE)=I^n,\qquad
 R^q p_*\OO_{X'}(-nE)=0\quad(q>0,\ n\text{ sufficiently large}).
\]
These are ordinary powers, not integral closures.
The relevant general facts are \cite[Tags 0AG6 and 0AG7]{Stacks}.
Projection, Leray, and Serre vanishing now give
\[
 H^1(X,I^n\otimes L^n)=
 H^1(X',\OO_{X'}(n(A-E)))=0
\]
for large \(n\).
Hence global sections of \(L^n\) realize every local class modulo
\(I^n\). Each generator of \(I\) has weight \(R\), so
\[
 I^n\subseteq\{\text{series of weight at least }nR\}.
\]
Surjectivity onto the former quotient therefore implies surjectivity
onto the desired smaller coefficient quotient.
Completion does not change the finite-length local quotients.
Finally, for sufficiently large \(n\), global sections of \(L^n\)
are restrictions of homogeneous polynomials of degree \(n\) in the
embedding coordinates, by Serre vanishing for the projective ideal
sheaf. They have \(W\)-degree at most \(nR\) on the affine chart.
Taking \(H=nR\) and undoing the triangular logarithmic substitution
proves \eqref{eq:jetmap}.
\hfill\(\square\)
